% 张三慧《大学物理学》（第4版）第二卷第14章
% 教材PDF物理页74--75，原印刷页57--58。
% 不要把页73--74的“思考题14.1--14.11”错当成章末“习题”。
\section*{第14章\quad 静电场中的导体}
\subsection*{章末习题14.1--14.14}

\paragraph{习题 14.1}
求导体外表面紧邻处场强的另一方法。
设导体面上某处面电荷密度为 \(\sigma\)，
在此处取一小面积 \(\Delta S\)，
将 \(\Delta S\) 两侧的电场看成是
\(\Delta S\) 面上的电荷的电场（用无限大平面算）
和导体上其他地方以及导体外的电荷的电场
（这电场在 \(\Delta S\) 附近可以认为是均匀的）的叠加，
并利用导体内合电场应为零求出导体表面处紧邻场强
为 \(\sigma/\varepsilon_0\)
（即式(14.2)）。

\paragraph{习题 14.2}
一导体球半径为 \(R_1\)，其同心地置于内、外半径
分别为 \(R_2,R_3\) 的厚导体壳。
此系统带电后内球电势为 \(\varphi_1\)，
外球所带总电量为 \(Q_2\)。
求此系统各处的电势和电场分布。
\textit{半径条件：}\(R_1<R_2<R_3\)，
以无穷远为电势零点。

\paragraph{习题 14.3}
在一半径为 \(R_1=6.0\,\mathrm{cm}\) 的金属球
\(A\) 外面套有一个同心的金属球壳 \(B\)。
已知球壳 \(B\) 的内、外半径分别为
\(R_2=8.0\,\mathrm{cm}\)，\(R_3=10.0\,\mathrm{cm}\)。
设 \(A\) 球带有总电量
\(Q_A=3\times10^{-8}\,\mathrm C\)，
球壳 \(B\) 带有总电量
\(Q_B=2\times10^{-8}\,\mathrm C\)。
（1）求球壳 \(B\) 内、外表面上各带有的电量
以及球 \(A\) 和球壳 \(B\) 的电势；
（2）将球壳 \(B\) 接地然后断开，
再把金属球 \(A\) 接地，
求金属球 \(A\) 和球壳 \(B\) 内、外表面上
各带有的电量以及金属球 \(A\) 和球壳 \(B\) 的电势。

\paragraph{习题 14.4}
一个接地的导体球，半径为 \(R\)，原来不带电。
今将一点电荷 \(q\) 放在球外距球心的距离为 \(r\) 的地方，
求球上的感生电荷总量（\(r>R\)）。

\paragraph{习题 14.5}
如图14.14所示，有三块互相平行的导体板，
外面的两块用导线连接，原来不带电。
中间一块上所带总面电荷密度为
\(1.3\times10^{-5}\,\mathrm{C/m^2}\)。
求每块板的两个表面的面电荷密度各是多少？
（忽略边缘效应）。
\textit{原图：}三板自上而下分别为导体A、B、C；
A、C以导线相连，A与B的间距为
\(5.0\,\mathrm{cm}\)，A与C的间距为
\(8.0\,\mathrm{cm}\)，因此B与C相距
\(3.0\,\mathrm{cm}\)。

\paragraph{习题 14.6}
一球形导体 \(A\) 含有两个球形空腔，
这导体本身的总电荷为零，
但在两空腔中心分别有一点电荷 \(q_b\) 和 \(q_c\)，
导体球外距导体球很远的 \(r\) 处有另一点电荷
\(q_d\)（图14.15）。
试求 \(q_b,q_c,q_d\) 各受到多大的力。
哪个答案是近似的？
\textit{图示：}两个点电荷位于各自独立球形空腔的球心，
\(r\) 从导体球几何中心量至远处 \(q_d\)。

\paragraph{习题 14.7}
试证静电平衡条件下导体表面单位面积受的力为
\[
 \boldsymbol f=\frac{\sigma^2}{2\varepsilon_0}\boldsymbol e_n,
\]
其中 \(\sigma\) 为面电荷密度，
\(\boldsymbol e_n\) 为表面的外法线方向的单位矢量。
此力方向与电荷的符号无关，
总指向导体外部。

\paragraph{习题 14.8}
在范德格拉夫静电加速器中，
是利用绝缘传送带向一个金属球壳输送电荷
而使球的电势升高的。
如果这金属球壳电势要求保持 \(9.15\,\mathrm{MV}\)，
（1）球周围气体的击穿强度为
\(100\,\mathrm{MV/m}\)，
这对球壳的半径有何限制？
（2）由于气体泄漏电荷，
要维持此电势不变，
需用传送带以 \(320\,\mu\mathrm{C/s}\)
的速率向球壳运送电荷。
这时所需最小功率多大？
（3）传送带宽 \(48.5\,\mathrm{cm}\)，
移动速率 \(33.0\,\mathrm{m/s}\)。
试求带上的面电荷密度和面上的电场强度。

\paragraph{习题 14.9}
一个点电荷 \(q\) 放在一无限大接地金属平板上方
\(h\) 处，
考虑到板面上紧邻处电场垂直于板面，
且板面上感生电荷产生的电场在板面上下具有对称性，
试根据电场叠加原理求出板面上感生面电荷密度的分布。

\paragraph{习题 14.10}
点电荷 \(q\) 位于一无限大接地金属板上方
\(h\) 处。
当问及将 \(q\) 移到无穷远需要做的功时，
第一个学生回答这功等于分开两个相距
\(2h\) 的电荷 \(q\) 和 \(-q\)
到无限远时做的功，即
\[
 A=\frac{q^2}{4\pi\varepsilon_0\cdot2h}.
\]
第二个学生求出 \(q\) 受的力再用
\(F\,dr\) 积分计算，从而得出了不同的结果。
第二个学生的结果是什么？
他们谁的结果对？

\paragraph{习题 14.11}
在图14.12中沿水平方向从 \(q\) 出发的那条电场线
在何处触及导体板？
（用高斯定律和简单积分。）

\paragraph{习题 14.12}
一条长直导线，均匀地带有电量
\(1.0\times10^{-8}\,\mathrm{C/m}\)，
平行于地面放置，且距地面
\(5.0\,\mathrm m\)。
导线正下方地面上的电场强度和面电荷密度各如何？
导线单位长度上受多大电力？

\paragraph{习题 14.13}
帕塞尔教授在他的《电磁学》中写道：
如果从地球上移去一滴水中的所有电子，
则地球的电势将会升高几百万伏。
请用数字计算证实他说的话。
\textit{数值建模：}原题没有指定一滴水的质量，
应明确列出按典型液滴质量估算的前提。

\paragraph{习题 14.14}
求半径为 \(R\)、带有总电量 \(q\)
的导体球的两半球之间的相互作用电力。
